% Copyright (c) 2026 by C. Klukas.
% Licensed under the MIT License. See LICENSE for details.

% This appendix is dense with monospaced names that cannot be hyphenated;
% the emergency stretch keeps such paragraphs inside the text width.
\begingroup\setlength{\emergencystretch}{3em}

The paglet ABI is the only contract between a host and a \paglet{} module.
It is plain core WebAssembly, without the component model: a few exports
the module provides, a few imports the host provides, and MessagePack
documents passed through linear memory. This appendix lists the contract as
the host implements it (\path{cpp/common/include/paglets/abi.hpp} holds the
constants and documents, \path{cpp/host/src/wasm_binary.cpp} the import
allowlists and the mutable-global check, \path{cpp/host/src/wasm_engine.cpp}
the import bindings and the pointer checks, \path{cpp/host/src/runtime.cpp}
the module check and the semantics); the design document \path{planning/cpp-abi-v1.md} is the
normative text and names the conformance tests. The guest SDK of
Section~\ref{sec:model} wraps all of this; a module written in another
language implements it directly.

\subsection{Module requirements and exports}\label{sec:abi-exports}

A module is a core Wasm module for \code{wasm32} with exactly one linear
memory, built as a reactor (\path{-mexec-model=reactor}). It declares the
ABI version by exporting a function \code{paglets\_abi\_v1}; the host never
calls it and looks only at the name. A module may declare several versions,
and the host uses the highest it supports; the current host supports v1
only and rejects a module that declares none. Table~\ref{tab:abi-exports}
lists the exports. The host checks them, the import set
(Section~\ref{sec:abi-imports}) and the mutable globals statically, before
it instantiates the module, and names the missing export or the offending
import in its error.

Every mutable global the module defines must be exported, because a memory
image captures the linear memory and the exported mutable globals and
nothing else (Section~\ref{sec:movement},
Listing~\ref{lst:movement-snapshot-ready}). A C or C++ guest built as
the documentation prescribes (a non-PIC reactor,
Appendix~\ref{app:reproduce}) has one such global, the stack pointer, so
the guest build links with \path{-Wl,--export=__stack_pointer}; the host
names any other unexported mutable global in its error. The module must
not change its tables at run time, so that function pointers stay valid
across images; C and C++ guests satisfy this, and the host cannot verify
it. The reactor initializer \code{\_initialize}, which runs static
constructors, takes effect once, at the \paglet{}'s creation: for a module
that imports WASI the runtime runs it at every instantiation and a restore
from an image then overwrites its effects; for a module without WASI
imports the host runs it at creation only.

\begin{table}[htbp]
\centering\small
\begin{tabular}{lp{42mm}p{58mm}}
\toprule
Export & Signature & Purpose \tabularnewline
\midrule
\code{paglets\_abi\_v1} & \code{() -> ()} & \raggedright version marker; never called \tabularnewline
\code{memory} & memory & \raggedright the one linear memory \tabularnewline
\code{paglets\_alloc} & \code{(size: i32) -> i32} & \raggedright a buffer for a document from the host; 0 on failure, which fails the \paglet{} like a trap \tabularnewline
\code{paglets\_free} & \code{(ptr: i32) -> ()} & \raggedright releases such a buffer \tabularnewline
\code{paglets\_on\_event} & \code{(kind, ptr, len) -> i32} & \raggedright lifecycle events (Table~\ref{tab:abi-events}) \tabularnewline
\code{paglets\_on\_message} & \code{(ptr, len) -> i32} & \raggedright message delivery (Table~\ref{tab:abi-delivered}) \tabularnewline
mutable globals & global & \raggedright every mutable global the module defines, for images \tabularnewline
\code{\_initialize} & \code{() -> ()} & \raggedright optional reactor initializer \tabularnewline
\code{paglets\_state\_export} & \code{() -> i64} & \raggedright reserved for code upgrades; accepted, never called \tabularnewline
\code{paglets\_state\_import} & \code{(ptr, len) -> i32} & \raggedright reserved for code upgrades; accepted, never called \tabularnewline
\bottomrule
\end{tabular}
\caption{The exports of an ABI v1 module. The first six are required (\code{abi::required\_exports}); the host rejects a module that lacks one. Unlabelled parameters are \code{i32}.}
\label{tab:abi-exports}
\end{table}

\subsection{Calling convention and memory ownership}\label{sec:abi-memory}

Figure~\ref{fig:abi-call} shows one delivery. To pass a document of
\code{len} bytes, the host calls \code{paglets\_alloc(len)}, copies the
document into the returned buffer, calls the handler with the pointer and
the length, and calls \code{paglets\_free} after the handler returns. The
guest must not keep pointers into the buffer. An empty document is passed
as \code{(0, 0)} without an allocation. In the other direction, imports
take pointer and length pairs; the runtime validates every pair against
the bounds of linear memory before the native function runs, and traps the
guest otherwise, so a guest cannot read host memory. The host copies what
it keeps before the import returns. Imports that return a document
(\code{self\_info}, \code{cap\_inspect}, \code{cap\_list}) write it into a
guest buffer \code{(out, cap)} and return the document's length; if the
length exceeds \code{cap}, nothing is written and the guest retries with a
larger buffer. A negative result is an error code. A document longer than
1~MB (\code{abi::max\_document\_size}) is refused with \code{too\_large}
before it is decoded.

\begin{figure}[htbp]
\centering
\begin{tikzpicture}[x=1cm,y=1cm]
  \node[box] (host) at (0,0) {host};
  \node[box] (guest) at (8,0) {\paglet{} module};
  \draw[lifeline] (host.south) -- ++(0,-6.6);
  \draw[lifeline] (guest.south) -- ++(0,-6.6);
  \draw[msg] (0,-0.8) -- node[above] {\code{paglets\_alloc(len)}} (8,-0.8);
  \draw[msg] (8,-1.5) -- node[above] {\code{ptr}} (0,-1.5);
  \node[event] at (4,-2.05) {copies the document to \code{ptr}};
  \draw[msg] (0,-2.9) -- node[above] {\code{paglets\_on\_message(ptr, len)}} (8,-2.9);
  \draw[msg] (8,-3.7) -- node[above] {\code{paglets.send}, \code{paglets.reply}, \ldots} (0,-3.7);
  \draw[msg] (0,-4.4) -- node[above] {\code{0} or an error code} (8,-4.4);
  \draw[msg] (8,-5.2) -- node[above] {\code{0} handled, \code{1} not handled, or an error} (0,-5.2);
  \draw[msg] (0,-6.0) -- node[above] {\code{paglets\_free(ptr)}} (8,-6.0);
  \node[event] at (4,-6.45) {the next call, or a memory image};
\end{tikzpicture}
\caption{One delivery. Imports run inside the handler; the host takes memory images only between two calls.}
\label{fig:abi-call}
\end{figure}

The host calls into a \paglet{} from one thread at a time and takes memory
images only at quiescent points, between two calls and never while an
import runs. \code{paglets\_on\_event} and \code{paglets\_on\_message}
return 0 for handled, 1 for not handled, or a negative error code
(Table~\ref{tab:abi-errors}). For a request whose handler returns 1 or an
error without having replied, the host answers the requester with status
\code{unknown\_message} or that error. A trap (\code{unreachable}, an
out-of-bounds access, stack exhaustion, the exceeded time budget of 5~s,
\code{proc\_exit}) fails the \paglet{}: the host discards the instance,
answers every pending request of the \paglet{} and every reply capability
it holds with \code{failed}, and ends it.

\subsection{Imports}\label{sec:abi-imports}

The host provides the module \code{paglets} with fourteen functions
(Table~\ref{tab:abi-imports}) and a subset of WASI preview~1
(\path{wasi_snapshot_preview1}). Pointers are \code{i32} offsets into
linear memory; \code{endpoint}, \code{handle} and \code{reply} are
capability handles, positive \code{i32} values indexing the \paglet{}'s
host-side capability table. Handle 1 is always the \paglet{}'s endpoint to
itself. Every use of a handle is checked against the table, so handles
cannot be forged; the table is host state and travels with the \paglet{}
as part of its record, not inside the image.

\begin{table}[htbp]
\centering\small
\begin{tabular}{lp{54mm}p{58mm}}
\toprule
Import & Signature & Result \tabularnewline
\midrule
\code{log} & \code{(level, ptr, len) -> ()} & \raggedright a log line; levels 0 debug to 3 error \tabularnewline
\code{self\_info} & \code{(out, cap) -> i32} & \raggedright length of the self document \tabularnewline
\code{send} & \code{(endpoint, msg, len) -> i32} & \raggedright 0 or an error \tabularnewline
\code{request} & \code{(endpoint, msg, len) -> i64} & \raggedright correlation ID ($>0$) or an error \tabularnewline
\code{reply} & \code{(reply, msg, len) -> i32} & \raggedright 0 or an error; consumes the reply capability \tabularnewline
\code{cap\_derive} & \code{(handle, spec, len) -> i32} & \raggedright a new, weaker handle or an error \tabularnewline
\code{cap\_drop} & \code{(handle) -> i32} & \raggedright 0 or an error; cancels a timer, answers a request \code{gone} \tabularnewline
\code{cap\_inspect} & \code{(handle, out, cap) -> i32} & \raggedright length of the capability document \tabularnewline
\code{cap\_list} & \code{(out, cap) -> i32} & \raggedright length of the list of held handles \tabularnewline
\code{create\_child} & \code{(spec, len) -> i32} & \raggedright endpoint handle to the new \paglet{}, or an error \tabularnewline
\code{lifecycle} & \code{(op, arg, len) -> i32} & \raggedright 0, or an endpoint handle for \code{clone}, or an error (Table~\ref{tab:abi-lifecycle}) \tabularnewline
\code{timer\_set} & \code{(delay: i64, msg, len) -> i32} & \raggedright timer handle or an error \tabularnewline
\code{now} & \code{(clock) -> i64} & \raggedright nanoseconds; clock 0 wall time since the Unix epoch, 1 monotonic \tabularnewline
\code{random} & \code{(out, len) -> i32} & \raggedright 0; fills the buffer with secure random bytes \tabularnewline
\bottomrule
\end{tabular}
\caption{The \code{paglets} import module (\code{abi::import\_names}). Unlabelled parameters are \code{i32}.}
\label{tab:abi-imports}
\end{table}

The WASI subset is the one of Table~\ref{tab:model-imports}:
\code{args\_get}, \code{args\_sizes\_get}, \code{environ\_get} and
\code{environ\_sizes\_get} (always empty); \code{clock\_res\_get},
\code{clock\_time\_get}, \code{random\_get} and \code{sched\_yield};
\code{proc\_exit}, which fails the \paglet{}; \code{fd\_write}, where
stdout and stderr go to the host log line by line at the levels info and
warning, and other descriptors are \code{EBADF}; and \code{fd\_close},
\code{fd\_seek}, \code{fd\_fdstat\_get}, \code{fd\_prestat\_get} and
\code{fd\_prestat\_dir\_name}, with no directory preopened. The host
implements \code{fd\_write} itself; the others come from the runtime's
WASI library. Files, sockets and every other resource are reached only by
message, through the system \paglet{}s of Section~\ref{sec:mesh}.

\subsubsection{Import policies}

Which imports a module may have depends on its trust class
(\code{wasm::ImportPolicy}). The \emph{standard} policy allows the
fourteen \code{paglets} functions and the WASI subset, and applies to
roaming and resident \paglet{}s. The \emph{system} policy adds
\code{paglets\_sys.*},
the module reserved for privileged host calls of system \paglet{}s written
as Wasm modules. In the current release that module is empty: the host
registers no functions under the name, and every system \paglet{} is
native. An import outside even the system set is refused when the module
is loaded. A roaming or resident module that imports from
\code{paglets\_sys} is accepted into the module store but refused when a
\paglet{} of that trust class is created from it, before instantiation,
with a message that names every offending import. The command
\code{paglets-host module inspect} prints a module's imports with the
policy's verdict and the trust classes it can run as.

\subsection{Events}\label{sec:abi-events}

\code{paglets\_on\_event(kind, ptr, len)} delivers the seven lifecycle
events of Table~\ref{tab:abi-events} (\code{abi::EventKind},
Listing~\ref{lst:lifecycle-abi-kinds}; the documents are
Listing~\ref{lst:lifecycle-abi-documents}). \code{created} and
\code{cloned} are the first delivery, ahead of any message already queued;
the others belong to their transition and are not queued. A negative
result is logged and does not stop the transition.
Appendix~\ref{app:events} gives the ordering of hooks and handlers.

\begin{table}[htbp]
\centering\small
\begin{tabular}{rlp{40mm}p{66mm}}
\toprule
Kind & Event & Document & Delivered \tabularnewline
\midrule
1 & \code{created} & \raggedright \code{StartEvent}: \code{args}, \code{caps} & \raggedright once, before any message \tabularnewline
2 & \code{activated} & none & \raggedright after a restore from its memory image (from the state directory or from memory), right before the delivery that woke the \paglet{} \tabularnewline
3 & \code{deactivating} & none & \raggedright after the handler that requested \code{deactivate}, before the image is taken \tabularnewline
4 & \code{arrived} & \raggedright \code{ArrivedEvent}: \code{from}, \code{lost} & \raggedright on the new host after a move; \code{lost} lists capabilities that could not be re-created there \tabularnewline
5 & \code{cloned} & \raggedright \code{StartEvent} with \code{original} & \raggedright in a clone, instead of \code{created} \tabularnewline
6 & \code{dispatching} & \raggedright \code{DispatchingEvent}: \code{destination} & \raggedright after the handler that requested \code{dispatch}, before the image is taken \tabularnewline
7 & \code{disposing} & none & \raggedright after the handler that requested \code{dispose}, before the end \tabularnewline
\bottomrule
\end{tabular}
\caption{Lifecycle events and their documents.}
\label{tab:abi-events}
\end{table}

\subsection{Messages and documents}\label{sec:abi-messages}

All documents are MessagePack maps with string keys; decoders skip keys
they do not know, which is what makes additions within a major version
compatible. Table~\ref{tab:abi-delivered} is the argument of
\code{paglets\_on\_message}; Table~\ref{tab:abi-outgoing} is the document
of \code{send}, \code{request}, \code{reply} and \code{timer\_set}.
Message names have 1 to 128 bytes, and names starting with
\code{paglets.} are reserved for the host: \code{send}, \code{request}
and \code{timer\_set} refuse them with \code{invalid\_argument}
(\code{reply} ignores the name; the delivered reply is always named
\code{reply}). Three such names carry the host's reports
(Table~\ref{tab:abi-reports}); each arrives as an ordinary message
(kind 0) at priority 7 and without a sender record; the SDK recognizes
them by their reserved name and the missing sender and turns them into
hooks (Listing~\ref{lst:lifecycle-sdk-reports}).

\begin{table}[htbp]
\centering\small
\begin{tabular}{llp{96mm}}
\toprule
Key & Type & Meaning \tabularnewline
\midrule
\code{kind} & int & \raggedright 0 message, 1 request, 2 reply, 3 timer (\code{abi::MessageKind}) \tabularnewline
\code{name} & str & \raggedright the message name \tabularnewline
\code{payload} & bin & \raggedright the payload, usually MessagePack \tabularnewline
\code{caps} & [int] & \raggedright transferred capabilities, already in the receiver's table \tabularnewline
\code{priority} & int & \raggedright 0 to 7 \tabularnewline
\code{sender} & map & \raggedright host-stamped record of the sender: \code{id}, \code{owner}, \code{trust}, \code{module}, \code{host}; absent for timers, replies and the host's reports \tabularnewline
\code{badge} & str & \raggedright badge of the endpoint the sender used, if any \tabularnewline
\code{reply} & int & \raggedright requests only: the one-shot reply capability \tabularnewline
\code{correlation} & int & \raggedright replies only: the ID that \code{request} returned \tabularnewline
\code{status} & int & \raggedright replies only: 0, or \code{timeout}, \code{unknown\_message}, \code{failed}, \code{gone}, or the error the handler returned \tabularnewline
\bottomrule
\end{tabular}
\caption{The delivered message (\code{abi::Delivered}).}
\label{tab:abi-delivered}
\end{table}

\begin{table}[htbp]
\centering\small
\begin{tabular}{llp{92mm}}
\toprule
Key & Type & Meaning \tabularnewline
\midrule
\code{name} & str & \raggedright required; \code{reply} for a reply \tabularnewline
\code{payload} & bin & \raggedright default empty \tabularnewline
\code{caps} & [int] & \raggedright handles to transfer; each must be transferable and moves to the receiver on success \tabularnewline
\code{priority} & int & \raggedright default 3; replies arrive at 7 \tabularnewline
\code{timeout\_ms} & int & \raggedright \code{request} only, positive; default 30\,000 \tabularnewline
\code{lend} & [int] & \raggedright v1.1: capabilities shown to a system \paglet{} for this message; they stay with the sender \tabularnewline
\code{async} & bool & \raggedright v1.1, \code{send} and \code{reply} only: return at once; a failure is reported as \code{paglets.undelivered} \tabularnewline
\bottomrule
\end{tabular}
\caption{The outgoing message (\code{abi::OutMessage}).}
\label{tab:abi-outgoing}
\end{table}

\begin{table}[htbp]
\centering\small
\begin{tabular}{lp{46mm}p{58mm}}
\toprule
Name & Payload & When \tabularnewline
\midrule
\code{paglets.undelivered} & \raggedright \code{Undelivered}: \code{name}, \code{handle}, \code{status} & \raggedright v1.1: an asynchronous \code{send} or \code{reply} could not be delivered; \code{status} is the error the synchronous call would have returned \tabularnewline
\code{paglets.move\_failed} & \raggedright \code{MoveFailed}: \code{destination}, \code{reason}, \code{clone} & \raggedright v1.1: a \code{dispatch} or a clone for another host did not happen; \code{clone} is the clone's ID, empty for a dispatch \tabularnewline
\code{paglets.recovered} & \raggedright \code{Recovered}: \code{image\_ms}, \code{restarted\_ms} & \raggedright v1.2: the host restarted and the \paglet{} came back from its last image; only with \code{notify\_recovered} \tabularnewline
\bottomrule
\end{tabular}
\caption{The host's reports: messages of kind 0 at priority 7 without a sender record.}
\label{tab:abi-reports}
\end{table}

Table~\ref{tab:abi-documents} lists the remaining documents. Capability
kinds are \code{endpoint}, \code{reply} and \code{timer}, and since v1.1
the resource kinds \code{dir}, \code{artifact}, \code{topic} and
\code{pin} handed out by system \paglet{}s; the design document also
names a \code{file} kind, which no service of the current release hands
out.

\begin{table}[htbp]
\centering\small
\begin{tabular}{lp{26mm}p{88mm}}
\toprule
Document & Import & Keys \tabularnewline
\midrule
self (\code{SelfInfo}) & \code{self\_info} & \raggedright \code{id}, \code{module} (hex SHA-256), \code{owner}, \code{trust}, \code{host}, \code{abi}, \code{minor}, \code{services} (service name to handle), \code{pinned\_until}, \code{restart}, \code{notify\_recovered} \tabularnewline
derive (\code{DeriveSpec}) & \code{cap\_derive} & \raggedright \code{ops}, \code{expires\_ms}, \code{uses}, \code{transferable}, \code{badge}, \code{path}; each optional, each only narrows the source \tabularnewline
capability (\code{CapInfo}) & \code{cap\_inspect} & \raggedright \code{kind}, \code{target}, \code{ops}, \code{transferable}, and when set \code{badge}, \code{expires}, \code{uses\_left}, \code{revoked} \tabularnewline
child (\code{ChildSpec}) & \code{create\_child} & \raggedright \code{module} (default: the creator's), \code{args}, \code{caps} \tabularnewline
\bottomrule
\end{tabular}
\caption{The documents of the capability and introspection imports (\code{abi.hpp}).}
\label{tab:abi-documents}
\end{table}

\subsection{Lifecycle operations}\label{sec:abi-lifecycle}

\code{lifecycle(op, arg, len)} requests a transition
(Table~\ref{tab:abi-lifecycle}). At most one of \code{dispose},
\code{deactivate} and \code{dispatch} can be pending per handler call;
they take effect after the handler returns, after any clones of the same
call (Section~\ref{sec:lifecycle-hooks}). An unknown operation returns
\code{invalid\_argument}.

\begin{table}[htbp]
\centering\small
\begin{tabular}{rlp{33mm}p{73mm}}
\toprule
Op & Name & Argument & Effect and errors \tabularnewline
\midrule
1 & \code{dispose} & none & \raggedright ends the \paglet{} after the handler; \code{bad\_state} if an end is pending \tabularnewline
2 & \code{deactivate} & \raggedright \code{wake\_after\_ms?} & \raggedright takes the image after the handler (to the state directory when the host has one, otherwise in memory); the next message, a timer or the delay activates it; \code{malformed} for an undecodable argument, \code{invalid\_argument} for a negative delay, \code{bad\_state} if an end is pending \tabularnewline
3 & \code{clone} & \raggedright \code{args}, \code{caps}, \code{destination?} & \raggedright a copy from the image taken after the handler; returns an endpoint to it at once; a system \paglet{}'s clone is roaming; \code{quota} beyond 16 per call or the capability limit; a local clone: \code{denied} when the policy refuses the module for the clone, \code{quota} at the owner's limit; with \code{destination} (v1.1): \code{unsupported} outside a mesh, \code{denied} unless roaming, \code{invalid\_argument} unless \code{caps} holds only endpoints to \paglet{}s and mesh artifacts \tabularnewline
4 & \code{dispatch} & \raggedright \code{destination} & \raggedright moves the \paglet{} after the handler; \code{unsupported} outside a mesh, \code{denied} for system and resident \paglet{}s, \code{pinned} while pinned, \code{bad\_state} if an end is pending \tabularnewline
5 & \code{checkpoint} & none & \raggedright v1.2: stores the image right after the handler, ahead of the host's periodic checkpoints (Section~\ref{sec:lifecycle-restart}); on a host with a state directory and unless the restart policy is off \tabularnewline
6 & \code{restart} & \raggedright \code{RestartArg}: \code{restart}, \code{notify\_recovered} & \raggedright v1.2: the restart policy (Listing~\ref{lst:lifecycle-restart}); \code{denied} for system \paglet{}s; clones inherit it \tabularnewline
\bottomrule
\end{tabular}
\caption{Lifecycle operations (\code{abi::LifecycleOp}).}
\label{tab:abi-lifecycle}
\end{table}

\subsection{Error codes}\label{sec:abi-errors}

Imports and handlers return the negative codes of
Table~\ref{tab:abi-errors} (\code{abi::Error}; \code{abi::error\_name}
gives the name). The same codes are reply statuses, the \code{status} of
an undelivered report, and the \code{error: <name>} of the command line.
An unknown code is a generic error to a v1.0 guest.

\begin{table}[htbp]
\centering\small
\begin{tabular}{rlp{104mm}}
\toprule
Code & Name & Meaning \tabularnewline
\midrule
$-1$ & \code{invalid\_argument} & \raggedright a parameter is out of range, or a document is incomplete \tabularnewline
$-2$ & \code{bad\_handle} & \raggedright the capability handle does not exist, or has the wrong kind \tabularnewline
$-3$ & \code{denied} & \raggedright the capability, or the mesh policy, does not allow the operation \tabularnewline
$-4$ & \code{expired} & \raggedright the capability has expired \tabularnewline
$-5$ & \code{quota} & \raggedright a quota or limit is exhausted: a full mailbox, too many \paglet{}s of an owner, too many pins, the storage quota, the uses of a capability \tabularnewline
$-6$ & \code{too\_large} & \raggedright a document or message exceeds its size limit \tabularnewline
$-7$ & \code{malformed} & \raggedright a body is not valid MessagePack, or does not decode into the expected type \tabularnewline
$-8$ & \code{not\_found} & \raggedright the target \paglet{}, module or service does not exist \tabularnewline
$-9$ & \code{bad\_state} & \raggedright not possible now, for example replying twice \tabularnewline
$-10$ & \code{unsupported} & \raggedright not supported by this host \tabularnewline
$-11$ & \code{internal} & \raggedright an error of the host \tabularnewline
$-12$ & \code{timeout} & \raggedright reply status: no reply within the request's timeout \tabularnewline
$-13$ & \code{unknown\_message} & \raggedright reply status: the receiver has no handler for the message \tabularnewline
$-14$ & \code{failed} & \raggedright reply status: the receiver trapped or exceeded the time budget; it has ended \tabularnewline
$-15$ & \code{gone} & \raggedright reply status: the receiver was disposed of, or did not answer the request \tabularnewline
$-16$ & \code{revoked} & \raggedright v1.1: the capability, one it was derived from, or its grant was revoked \tabularnewline
$-17$ & \code{pinned} & \raggedright v1.1: \code{dispatch} while the \paglet{} is pinned to its host \tabularnewline
\bottomrule
\end{tabular}
\caption{Error codes.}
\label{tab:abi-errors}
\end{table}

The limits behind \code{quota} and \code{too\_large} are compiled-in
defaults of \code{runtime::Config} and \code{wasm::Limits}: 1~MB per
document, 1024 messages per mailbox, 1024 capabilities and 64 pending
timers per \paglet{}, 16 children and clones per handler call, 5~s per
handler call, 16~MB of linear memory, 10\,000 \paglet{}s of one owner per
host and 64 pins on one \paglet{}.

\subsection{Version history}\label{sec:abi-versions}

Within the major version 1, additions are compatible: a module using an
import the host does not know is rejected at load time, both sides ignore
document keys they do not know, and a guest treats unknown event kinds as
no-ops and unknown error codes as generic errors (the SDK returns
\code{not\_handled} for an unknown kind, \code{abi::error\_name} returns
\code{error} for an unknown code). The host reports its
minor version in \code{self\_info.minor} (\code{abi::minor\_version}, 2 in
the current release), so a guest can test for a feature before it uses
it. Table~\ref{tab:abi-versions} lists the additions.

\begin{table}[htbp]
\centering\small
\begin{tabular}{lp{124mm}}
\toprule
Version & Content \tabularnewline
\midrule
v1.0 & \raggedright the fourteen imports, the six required exports, the seven event kinds, the four message kinds, the operations \code{dispose}, \code{deactivate}, \code{clone} and \code{dispatch} (the last answering \code{unsupported}), the capability kinds \code{endpoint}, \code{reply} and \code{timer}, the error codes $-1$ to $-15$ \tabularnewline
v1.1 & \raggedright resource capabilities (\code{dir}, \code{artifact}, \code{topic}, \code{pin}; the design document also names \code{file}) with \code{revoked} ($-16$) and the derive \code{path}; \code{lend}; default service endpoints in \code{self\_info.services}; \code{async} with \code{paglets.undelivered}; \code{dispatch} and \code{clone} to another host in a mesh, with \code{arrived}, \code{dispatching} and \code{paglets.move\_failed}; \code{pinned} ($-17$) and \code{pinned\_until}; \code{minor} \tabularnewline
v1.2 & \raggedright \code{checkpoint} and \code{restart} with \code{RestartArg}; \code{paglets.recovered}; \code{restart} and \code{notify\_recovered} in \code{self\_info} \tabularnewline
\bottomrule
\end{tabular}
\caption{Additions within ABI v1.}
\label{tab:abi-versions}
\end{table}

The design document lists 41 conformance tests, C01 to C41, from the
rejection of a module without the version marker to the behaviour of a
pinned \paglet{}. They run in CI against the host runtime: C01 to C30 and
C36 in \path{cpp/tests/test_runtime.cpp} with the conformance guest
\path{cpp/tests/guests/conformance.cpp}, the system-\paglet{} tests C31
to C35 in \path{cpp/tests/test_system.cpp}, the movement tests C37 to C40
in \path{cpp/tests/test_mobility.cpp} and the pin test C41 (unnumbered
there) in \path{cpp/tests/test_location.cpp}.
\par\endgroup
